 \documentclass{article}
\usepackage{tikz}
\usetikzlibrary{angles, quotes}
\usepackage{amssymb}
\usepackage{pgfplots}
\pgfplotsset{width=10cm,compat=1.9}
\usepgfplotslibrary{statistics}
\usepackage[left=2.54cm,right=2.54cm,top=2.54cm,bottom=2.54cm]{geometry}
\usepackage{tasks} %Paquete necesario para la producción de items horizontales para algunos ejercicios de tarea empleando el comando /begin{multicols}}{#}.../end{multicols} para ayudarnos a reducir las páginas
\usepackage{pdflscape} %comando que permite cambiar a orientación horizontal las hojas%
\usepackage{soul} %Paquete para permitir la compilación de acentos usando el comando {\hl{}}}\\
\sethlcolor{yellow(munsell)} %comando para definir el color para el subrayado usando el comando \hl
\usepackage[utf8]{inputenc}
\usepackage[spanish]{babel}
\usepackage{icomma}
\usepackage{ragged2e}
\usepackage{siunitx}
\usepackage{url}
\usepackage[colorlinks=true, urlcolor=blue,  linkcolor=black, citecolor=green]{hyperref}
\usepackage{pdfpages}
\usepackage{blindtext} %Paquete que genera texto ficticio (dummy text) con el comando: \blindtext
\usepackage{cancel} %in the preamble gives you four different modes of striking through: \cancel{text to cancel} draws a diagonal line (slash) through its argument, \bcancel{text to cancel} uses the negative slope (a backslash), \xcancel{text to cancel} draws an X (actually \cancel plus \bcancel), \cancelto{〈value〉}{〈expression〉} draws a diagonal arrow through the 〈expression〉pointing to the 〈value〉 (math-mode only)
\usepackage{csquotes}
\usepackage{afterpage}
\usepackage{parskip} 
\usepackage{float}
\usepackage{enumitem}
\usepackage{multicol}%Paquete que permite la creación aislada de columnas de texto. De esta forma se reduce la cantidad de páginas en nuestro documento
\usepackage{lipsum} %Paquete que genera texto de relleno al igual que \usepackage{blindtext}, se usa con el comando \lipsum[#-#] los #(hashtags) delimitan el número de páginas que deseamos ocupar del paquete lipsum para usarlos como relleno. 

% Definir el comando personalizado para la nota
\newcommand{\mynote}[1]{%
\begin{tikzpicture}[baseline=-0.75ex]
\node[draw, circle, fill=Apple Green!60, inner sep=2pt] (note) {\textbf{Nota:}};
\end{tikzpicture}%
\ \textit{#1}%
}

\newenvironment{Figura}
{\par\medskip\noindent\minipage{\linewidth}}
{\endminipage\par\medskip}
\usepackage{caption}
\usepackage[
backend=biber,
sorting=none,
url=true
]{biblatex} %bibliografía
\addbibresource{biblio.bib}
% Establecer el espacio entre las entradas de la bibliografía
\setlength{\bibitemsep}{1\baselineskip} % Puedes ajustar el valor según tus preferencias
\usepackage{amsmath,amsthm,amssymb,amsfonts}
\usepackage{pifont} %Permite la compilación del comando \xmark: es el símbolo de la tachita
\usepackage{mathtools} %permite la compilación de símbolos para matrices
\usepackage{empheq} %Paquete que se relaciona con \usepackage[most]{tcolorbox} para la creación de cuadros/cajas de colores para delimitar los resultados matemáticos o líneas de texto usando la línea de comando: \begin{empheq}[box={\mymath[colback=orange(webcolor)!70,drop lifted shadow]}]{equation*} \end{empheq}
\usepackage[most]{tcolorbox} %Paquete que se relaciona con \usepackage{empheq} para la creación de cuadros/cajas de colores para delimitar los resultados matemáticos o líneas de texto
\renewcommand{\qedsymbol}{$\blacksquare$}
\usetikzlibrary{positioning,decorations.pathreplacing} %Paquete que permite la compilación de llaves para cuadros sinópticos (brace diagram) 
\usepackage{schemata} %The schemata package is designed just for these brace diagrams
\usepackage{booktabs}

\newcommand\AB[2]{\schema{\schemabox{#1}}{\schemabox{#2}}} %comando o paquete necesario para crear cuadros sinópticos

\newtcbox{\mymath}[1][]{%
nobeforeafter, math upper, tcbox raise base,
enhanced, colframe=blue!30!black,
colback=blue!30, boxrule=1pt,
#1} %Comando importante para encasillar los resultados matemáticos en cajas de diferentes colores, se relaciona con los paquetes \usepackage{empheq} y \usepackage[most]{tcolorbox} 

\newenvironment{sysmatrix}[1]
{\left(\begin{array}{@{}#1@{}}}
{\end{array}\right)}
\newcommand{\ro}[1]{%
\xrightarrow{\mathmakebox[\rowidth]{#1}}%
}
\newlength{\rowidth}% row operation width
\AtBeginDocument{\setlength{\rowidth}{3em}}  %Comando importante para laproducción de lineas de operaciones entre matrices, método de Gauss o Gauss Jordan se relaciona con \newenvironment{sysmatrix}

%formato para cambiar el horario a español
\usepackage[spanish]{datetime2}
\DTMsetdatestyle{spanish}

\renewcommand{\today}{\DTMdisplaydate{\the\year}{\the\month}{\the\day }{-1}}
%%%%%%%%%%%%%%%%%%%%%%%%%%%%%%%%%%%%%%%%%

\usepackage{datetime}
\newcommand{\mycurrenttime}{\xxivtime}

%%%%%%%%%%%%%%%%%%%%%%%%%%%%%%%%%%%%%%%%%%%%%%%%%%%%%%%%%%%%%%%%%%%%%%%%%%%%%%%%%%%%%%%%%%%%%%%%%%%%%%%%%%%%%%%%%%%%%%%%%%%%%%%%%%%%%%%%%%%%%%%%%


%%%%%%%%%%%%%Caja de color para el título%%%%%%%%%%%%%%%%%%%%%%%%%%%%%%%%%%%%%%%%%%%%%%%%%%%%%%

\definecolor{myframecolor}{RGB}{1, 45, 75} % Prussian Blue
\definecolor{myboxcolor}{RGB}{0, 107, 92} % Apple Green

\newtcolorbox{mybox}{
enhanced,
colback=myboxcolor!7, % Color de fondo del cuadro
colframe=myframecolor, % Color del marco del cuadro
arc=0pt,
boxrule=1pt,
borderline west={2mm}{-10mm}{myframecolor}, % Borde en el lado izquierdo
sharp corners=southwest,
width=\linewidth,
before=\par\vspace{\bigskipamount}, % Espacio antes del cuadro
after=\par\vspace{\bigskipamount} % Espacio después del cuadro
}


\newcommand{\euler}{e} %Comando para producir letra e de euler
\newenvironment{remark}{\par\vfill\footnotesize % Vertical white space above the remark and smaller font size
\begin{list}{}{
	\leftmargin=80pt % Indentation on the left
	\rightmargin=60pt}\item\ignorespaces % Indentation on the right
\makebox[-2.5pt]{\begin{tikzpicture}[overlay]
		\node[draw=Horizon!60,line width=2.5pt,circle,fill=Horizon!25,font=\sffamily\bfseries,inner sep=4pt,outer sep=0pt] at (-19pt,5pt){\textcolor{Horizon}{Nota}};\end{tikzpicture}} % Blue Nota in a circle
\advance\baselineskip -1pt}{\end{list}\vskip5pt} % Tighter line spacing and white space after remark
\usepackage{graphicx}

\usepackage{titling}

\input{colores.tex}

%%Producir signos de cita textual``Asignatura''

\newcommand{\xmark}{\ding{55}} %%%comando que genera la tachita

\newenvironment{MyColorPar}[1]{%
\leavevmode\color{#1}\ignorespaces%
}{%
}%

%%%%%%%%%%%%%%%%%%% Título de la tarea, Nombre de alumno

\title{Clase \#35 - Electrodinámica}
\author{\emph{} $\boldsymbol{\mid}$}
\date{\today}

%%%%%%%%%%%%%%%%%%% Datos de la Materia

\usepackage{fancyhdr}
\fancypagestyle{plain}{%  the preset of fancyhdr 
\fancyhf{} % clear all header and footer fields
\fancyfoot[R]{
\includegraphics[width=2cm]{LogoFCFMBUAP (1).png}}
\fancyfoot[L]{{\bfseries{\thedate{}}} a las {\bfseries{\mycurrenttime{}}} horas{} (GMT-6, H. Puebla de Zaragoza, Pue)}
\fancyhead[L]{Electrodinámica}
\fancyhead[R]{\theauthor}
}
\makeatletter
\def\@maketitle{%
\newpage
\null
\vskip 1em%
\begin{center}%
\let \footnote \thanks
{\LARGE \@title \par}%
\vskip 1em%
%{\large \@date}%
\end{center}%
\par
\vskip 1em}
\makeatother

\usepackage{lipsum}  
\usepackage{cmbright}


%%%%%%%%%%%%%%%%%%%%%%%%%%%%%%%%%%%%%%%%%%%%%%%%%%%%%%%%%%%%%%%%%%%%%%%%%%%%%%%%%%%%%%%%%%%%%%%%%%%%%%%%%%%%%%%%%%%%%%%%%%%%%%%%%%%%%%%%%%%%%%%%%%%%%%%%%%%%%%%%%%%%%%%%%%%%%%%%%%%%%%%%%%%%%%%%%%%%%%

\begin{document}



\maketitle


%%%%%%%%%%%%%%%%%%% Datos del profesor, campus y aula

\begin{mybox}
\noindent\begin{tabular}{@{}ll}
	{\bfseries{Profesor}} & Dr. Iván Fuentecilla Cárcamo\\
	\textcolor{prussianblue}{\bfseries{Campus}} / \textcolor{trueblue}{\bfseries{Facultad}}  & \textcolor{prussianblue}{\bfseries{C.U. BUAP}} /  \textcolor{trueblue}{\bfseries{FCFM}} \\
	\textcolor{prussianblue}{\bfseries{Edificio}} /  \textcolor{trueblue}{\bfseries{Salón}}    & 	\textcolor{prussianblue}{\bfseries{EMA4}} /  \textcolor{trueblue}{\bfseries{409}}   
\end{tabular} 
\end{mybox} \vspace{1cm}




	





\begin{MyColorPar}{orange(colorwheel)}

\end{MyColorPar}


\tableofcontents \newpage

\begin{center}
\begin{tikzpicture}
	% Dibuja círculos rellenos en diferentes colores
	\fill[bittersweet] (0,0) circle (0.17cm);
	\fill[Apple Green] (0.6,0) circle (0.17cm);
	\fill[trueblue] (1.2,0) circle (0.17cm);
	\fill[Aguamarina] (1.8,0) circle (0.17cm);
	\fill[Sun] (2.4,0) circle (0.17cm);
	\fill[orange(colorwheel)] (3.0,0) circle (0.17cm);
	\fill[Pesto] (3.6,0) circle (0.17cm);
	\fill[blue-violet] (4.2,0) circle (0.17cm);
	\fill[ticklemepink] (4.8,0) circle (0.17cm);
	\fill[Gallery] (5.4,0) circle (0.17cm);
	\fill[Mantle] (6.0,0) circle (0.17cm);
	\fill[Nevada] (6.6,0) circle (0.17cm);
	\fill[onyx] (7.2,0) circle (0.17cm);
	% Agrega más círculos o modifica los colores según sea necesario
\end{tikzpicture}
\end{center}  \vspace{2cm}




\section{Unit 3: Magnetostatics}

\begin{figure}[H]
    \centering
	    \includegraphics[width=0.7\textwidth]{1.png} % ajusta el ancho según convenga
    \caption{Descripción clara y concisa de la figura.}
    \label{fig:mi_figura2}
\end{figure}





\textbf{Operador de velocidad} asociado a cada partícula cargada. Queremos obtener un promedio de todas las partículas cargadas, en escencia esas partículas cargadas, son partículas cuánticas.


De tal forma que cuando hablamos de estas partpiculas cargadas, podemos hablar de una cantidad promedio.

Tenemos:


\[ \langle  \hat{v}_{i} \rangle \quad \vec{drif\epsilon}  \]


Las partículas que conforman a la corriente eléctrica se mueven a la velocidad común. NO es es eso, más bien se resuelve la ecuación de Shrodinger, es decir se resuelve para un examen. Se obtiene una velocidad de deriva que se relaciona con lo que vemos a continuación.


\textbf{Promedio de valores esperados} para el ensamble $\langle \hat{v}_{i} \rangle$

Según Reitz-Mildford atraviesa una cantidad de carga sobre la supeerficie. Si $\Delta{t}$ se va reduciendo. 


Tenemos:


\[  \vec{v} \cdot \hat{n} \delta{\epsilon} da  \]

\textit{La idea es cuantificar cuanta carga tenemos de ese volumen que va a a travesar esa área $A$}.

\[\mathcal{N} = \text{Density of charge carriers}\]


Tomamos el \textbf{límite}:


\[  \lim\limits_{\delta \epsilon \to 0} \frac{  \mathcal{N}q\vec{v}\cdot \hat{n}\delta\epsilon da   }{\delta \epsilon} = I  \]

Con:

\[ I = Nq\vec{v}\cdot \hat{n}da \]

\begin{figure}[H]
    \centering
	    \includegraphics[width=0.7\textwidth]{2.png} % ajusta el ancho según convenga
    \caption{Descripción clara y concisa de la figura.}
    \label{fig:mi_figura2}
\end{figure}


\textbf{Densidad de corriente}

\[ \underbrace{\vec{J}}_{\text{Densidad de corriente}} = 	\sum\limits_{i}^{n} N_{i}q_{i}\vec{v}_{i}   \]

Nos da una direccionalidad y cuanta carga pasa por unidad de tiempo.

POdemos fijarnos en una superficie finita, no en un elemento diferencial de área.


Tenemos nosotros (\hl{ejemplo}) :

\textbf{Tazón (superficie finita, no infinitesimal)}

Tenemos un campo vectorial, para cada punto en el espacio tenemos un vector densidad de corriente.

\textbf{Pregunta:} \textit{¿La densidad corriente de masa dependerá del tiempo?}


\begin{figure}[H]
    \centering
	    \includegraphics[width=0.7\textwidth]{3.png} % ajusta el ancho según convenga
    \caption{Descripción clara y concisa de la figura.}
    \label{fig:mi_figura2}
\end{figure}

Así $J$ depende de $\vec{r}$ y del tiempo $t$. EN el caso de magnétostatica solo nos interesa la dependecia en $r$.


La densidad de corriente es un campo vectorial, en el tazón tendremos muchas direcciones.

Lo que calculamos es la corriente eléctrica en el diferencial de área, hacemos una integral para conocer en toda la superficie finita. 

SOn contribuciones en ese elemento diferencial de área.


Tenemos nosotros:

\textbf{Corriente asociada a esa superficie finita}

\[I= \int\limits_{S} \vec{J} \cdot \hat{n} da  \]

Con este resuktado anterior veremos lo que es la ecuación de continuidada.


EN los procesos nucleares existe conservacipon de carga, en un decaimiento mismo. i.e en la \textbf{aniquilación} antielectrón-electrón. 


\[  e^{+} + e^{-} \to \gamma \]  


Imaginemos que el tazón cierra sobre si mismo. Tenemos carga que cambia, tenemos un volumen con cierta cantidad de carga a cierto tiempm $t$. Tenemos que tener una ecuación que nos relacione esa conservación. A continuación leemos esta apartado en el Reitz-Mildford. 


\textbf{Convención}.


Acá estamos suponiendo que la corriente $i$ es positiva cuando tenemos flujo a fuera.

\begin{figure}[H]
    \centering
	    \includegraphics[width=0.7\textwidth]{4.png} % ajusta el ancho según convenga
    \caption{Descripción clara y concisa de la figura.}
    \label{fig:mi_figura2}
\end{figure}


La corriente es positiva cuando consideramos la corriente que circula de afuera hacia adentro.


\[ I=-\oint \vec{J} \cdot \hat{n}da \]


\textbf{Use the divergence theorem}:

\[ I = \int\limits_{v} \frac{\partial }{\partial t}  \rho(\vec{r},t) dv \]


\textbf{Obtener $I$ i n terms of $\rho(\vec{r},t)$}

\[ I = \int\limits_{V} \frac{\partial }{\partial t}   \rho (\vec{r},t)dv \]

\textbf{4) Use 2) and 3) to deduce continuity equation}:


\[  \int\limits_{V}  \frac{ \partial P }{\partial t} dV = - \int\limits_{V} \nabla \cdot \vec{J} dv  \]

\[  \int\limits_{V} \left(   \frac{ \partial P }{\partial t} + \nabla \cdot \vec{J}    \right) dv =0  \]


\[  \boxed{ \frac{\partial P}{\partial t}  + \nabla \cdot \vec{J} =0 } \]


\textcolor{Cinnabar}{\textbf{Sink}}:


\[ \boxed{\frac{\partial }{\partial t } = -\nabla \cdot \vec{J}} \]


Usando la \textbf{ley de Ohm} para el caso anterior:

Tenemos:

\underline{Conductividad} resulta ser al mismo nivel que la permitividad. 


\textbf{Ohm$^{\prime}$s law}


Tenemos:


\[ \boxed{ \vec{J} = \sigma \vec{E}  } \]


En general la $\sigma$ va a depender del campo eléctrico en general.


\[  \sigma = \sigma(\vec{E})  \]

and it is a \textbf{tensor}. 

Los tensores son importantes dado que, por ejemplo de forma natural elegimos el plano $xy$ como perpendicular a la dirección de la plomada. 


Si queremos realmente describir las leyes de la física, necesitamos da una dependencia de las leyes de la física. 

En ese sentido, necesitamos leyes físicas que sean invariantes ante transformaciones, es ahí donde surgen los tensores.

La direccionalidad se absorbe en la naturaleza tensorial. Resulta entonces que la $\sigma$ es una cantidad tensorial, pero para este curso la $\sigma$ es una constante. 


Pensemos en un modelo: NIños atados a sus papás, existe un campo eléctrico que hace que los niños se muevan, si por ejemplo tenemos un campo electromagnético, sería un campo electrico y magnético dependeientes del tiempo. Para la luz ambos se mueven perpendiculares, en el caso de la luz. Resulta que lo niños reaccionan al campo eléctrico. 


El campo eléctrico hace que el niño con menos masa pueda despazarse mejor, ahora la radiación electromagnética. 


A las ondas de radio, nos permite tener comunicaciones. La frecuencia es cuantas oscilaciones tiene el campo por uniad de tiempo.


Existirá una tendencia a moverse de los niños, si la frecuencia es muy alta, ¿qué pasará con esos niños? 


A los portadores de carga ligados, sucede que si la frecuencia es alta oscilan demasiado, mientras que si la frecuencia es suave, entonces se mueven de manera suave. 


Si entendimos la analogía, para el caso de una onda electromagnética, un campo eléctrico y magnético.


Para este caso:


\[ \vec{E}, H(\vec{r},t) \]


\[ \underbrace{\sigma(\omega)}_{\text{Optical conductivity}} \]

La sigma en efecto depende de la frecuencia, cuando tenemos eso, decimos que la conductivad es una conductividad óptica:


Definimos ahora a la \textbf{resistividad}.


\textbf{Resistivity}:


\[  \rho = \frac{1}{\sigma}  \]

La resistividad es una propiedad de los materiales, en el Reitz-Mildford aparece una tabla para los ordenes típicos e algunos metales. 


Entoces como se muestra en la tabla, la \textbf{conductividad} depende de la temperatura.


Ahora con todas estas ideas tenemos:


\textbf{Siguiente actividad}:


Tenemos que deducir la \textbf{Ley de Ohm} en la forma $\Delta{V} = RI$ utlizando la \textbf{ley de ohm} para materiales lineales (ohmic materials)

Donde:

\[R=\text{Resistance}\]

\[ I = \text{Corriente} \]



Tenemos:


\textbf{1) Integrate $\vec{J}$ for a linear wire}:


\[ I = \int\limits \vec{J} \cdot \hat{n} dA = JA  \]




\textbf{2) Use Ohm$^{\prime}$s law ($ \vec{J} = \sigma \vec{E} $) to relate $I$, $\Delta{V}$ and \underline{resistenance}}.



Tenemos:



\begin{equation}
\begin{split}
\vec{J}  &= \sigma \vec{E} \quad \to \quad J=\sigma E\\
I&=JA \quad \to \quad I=\sigma EA\\
\Delta{\varphi} &= \int\limits \vec{E} \cdot d\vec{\ell}= E\ell \quad \to \quad E = \frac{ \Delta{\varphi} }{\ell}\\
\to I &= \sigma \frac{\Delta{\varphi}}{\ell} A \quad \text{y} \quad R=\frac{\ell}{\sigma A}\\
\to \quad I&= \frac{I}{R}\Delta{\varphi} \quad \to \quad IR=\Delta{\varphi}
\end{split}
\end{equation}


\textbf{Comment}:

\[ \sigma(\omega) \]


\[ \vec{E} = \vec{A} e^{i \left[  \vec{k}\cdot \vec{r} - \omega{t}   \right]  } \]

\underline{Plane wave - onda plana.}

\[ \vec{J} = \sigma \vec{E} = \sigma (\omega) \vec{E}  \]

\[\vec{J} = \vec{J}( \vec{r},t,\omega )  \]



Esto lo implementamos para un material intresante, esto es el proyecto final. 

Tenemos que hallar el campo eléctrico en una región del espacio, tenemos también que calcular la densidad de energía.




\end{document}